\documentclass[10pt]{article}
\usepackage[a4paper,margin=20mm]{geometry}
\usepackage{amsmath,amssymb,amsthm,mathtools}
\usepackage{graphicx}
\usepackage{microtype}
\usepackage{xcolor}
\usepackage[colorlinks=true,linkcolor=blue!55!black,citecolor=blue!55!black,urlcolor=blue!55!black]{hyperref}

\newtheorem{theorem}{Theorem}
\newtheorem{lemma}[theorem]{Lemma}
\newtheorem{corollary}[theorem]{Corollary}
\newcommand{\D}{\mathbb D}
\newcommand{\Hbeta}{\mathcal H^2(\beta)}
\newcommand{\Lin}{\mathcal L}
\newcommand{\Dom}{\mathcal D}
\newcommand{\C}{\mathbb C}

\title{The Hyperbolic Composition Generator Is Never Bounded\\
\large and never has an $H^\infty$ calculus on a proper sector}
\author{Candidate full result for the Section 3 question in arXiv:1708.02259}
\date{13 August 2026}

\begin{document}
\maketitle

\begin{center}
\fbox{\parbox{0.92\linewidth}{\textbf{Status: candidate full proof, likely valid.}
The literal source question has a negative answer for every positive weight
sequence.  Under the natural alternative reading suggested by the surrounding
discussion, every admissible weight sequence gives the requested failure of a
proper-sector calculus.  Novelty remains subject to expert literature review.}}
\end{center}

\section{Source question and the two readings}

Let $\beta=(\beta_n)_{n\geq0}$ be positive and let
\[
 \Hbeta=\left\{f(z)=\sum_{n\geq0}a_nz^n:
       \sum_{n\geq0}|a_n|^2\beta_n^2<\infty\right\},
 \qquad \|z^n\|=\beta_n.
\]
As in the source, assume that this is a Hilbert space of holomorphic functions
satisfying
$\mathcal O(\overline\D)\hookrightarrow\Hbeta\hookrightarrow
\operatorname{Hol}(\D)$ continuously.  Put
\[
 \alpha_r(z)=\frac{z+r}{1+rz},\qquad
 \phi_t=\alpha_{\tanh t},\qquad T(t)=C_{\phi_t}.
\]
The displayed question on source PDF p.~7 is reproduced below \cite{GY}.

\begin{center}
\includegraphics[width=0.98\linewidth]{figures/source_question_crop.png}
\end{center}

\noindent\textbf{Exact transcription.}
``Do there exist weights $\{\beta_n\}$ such that the generator of the
semigroup $\{C_{\alpha_{\tanh t}}\}$ on $\mathcal H^2(\beta)$ is bounded,
but does not admit an $H^\infty$ calculus in a sector
$|\arg(z_0-z)|<\theta$, where $\theta\in(0,\pi/2)$?  The same can be asked
for an arbitrary bounded composition semigroup on $\mathcal H^2(\beta)$.''

Taken literally, the generator is required to be a bounded operator.  The
surrounding paragraphs instead contrast bounded composition operators with
quasicontractivity, so a plausible intended reading is that the
\emph{semigroup consists of bounded operators} while its generator lacks the
proper-sector calculus.  Theorem~\ref{thm:main} resolves both readings.
Here ``bounded composition semigroup'' is used in the source's pointwise
sense: each $T(t)$ is a bounded operator, not necessarily
$\sup_{t\geq0}\|T(t)\|<\infty$.

\subsection*{Proof Intuition}

The weight sequence can make one of the two shifts in
$(1-z^2)d/dz$ small, but not both.  On two adjacent normalized monomials,
the relevant shift coefficients are reciprocal up to the factor $n(n+1)$;
their product therefore diverges independently of the weights.  This rules
out a bounded generator.

For the functional-calculus statement, the hyperbolic maps form a group, not
merely a one-sided flow.  Reflection in the origin turns every positive-time
composition operator into its negative-time counterpart.  A proper-sector
$H^\infty$ calculus would make the semigroup analytic.  But an analytic
semigroup that also extends to a $C_0$-group has bounded generator, because
positive time maps the whole space into the generator domain and negative
time inverts that map.  The monomial obstruction then gives the contradiction.

\section{The weight-independent monomial obstruction}

The derivative of the flow at zero is
\[
 \left.\frac{d}{dt}\phi_t(z)\right|_{t=0}=1-z^2.
\]
Under the hypotheses above, the source's generator-domain formula gives
\begin{equation}\label{eq:generator}
 Af=(1-z^2)f',\qquad
 \Dom(A)=\{f\in\Hbeta:(1-z^2)f'\in\Hbeta\}.
\end{equation}
In particular every polynomial belongs to $\Dom(A)$.

\begin{lemma}[Adjacent-monomial obstruction]\label{lem:unbounded}
For every positive sequence $\beta$, the operator in
\eqref{eq:generator} is unbounded on $\Hbeta$.
\end{lemma}

\begin{proof}
Let $e_n=z^n/\beta_n$.  These are unit vectors, and for $n\geq1$,
\begin{equation}\label{eq:matrix}
 Ae_n=n\frac{\beta_{n-1}}{\beta_n}e_{n-1}
      -n\frac{\beta_{n+1}}{\beta_n}e_{n+1}.
\end{equation}
Orthogonality gives
\[
 \|Ae_n\|\geq n\frac{\beta_{n+1}}{\beta_n},\qquad
 \|Ae_{n+1}\|\geq(n+1)\frac{\beta_n}{\beta_{n+1}}.
\]
The product of the right-hand sides is $n(n+1)$.  Hence
\begin{equation}\label{eq:lower}
 \max\{\|Ae_n\|,\|Ae_{n+1}\|\}\geq\sqrt{n(n+1)}
 \longrightarrow\infty.
\end{equation}
Thus $A$ is unbounded on its domain, for every choice of positive weights.
\end{proof}

The use of adjacent inputs is essential.  Looking only at $Ae_n$ permits the
ratio $\beta_{n+1}/\beta_n$ to suppress its forward coefficient; the backward
coefficient on $Ae_{n+1}$ is then forced to grow by the reciprocal amount.

\section{The group and proper-sector obstruction}

\begin{lemma}[The positive flow automatically extends to a group]
Suppose $T(t)=C_{\phi_t}$, $t\geq0$, is a $C_0$-semigroup on $\Hbeta$.
Then the same formula for $t\in\mathbb R$ defines a $C_0$-group.
\end{lemma}

\begin{proof}
Let $\rho(z)=-z$ and $R=C_\rho$.  Since the monomials are orthogonal, $R$ is
unitary on $\Hbeta$.  Direct calculation gives
\[
 \rho\circ\alpha_r\circ\rho=\alpha_{-r},
 \qquad C_{\alpha_{-r}}=R C_{\alpha_r}R.
\]
Thus all negative-time operators are bounded and are conjugates by the fixed
unitary $R$.  Strong continuity for negative times follows from strong
continuity for positive times, and
$\phi_{t+s}=\phi_t\circ\phi_s$ gives the group law.
\end{proof}

\begin{lemma}[An analytic $C_0$-group has bounded generator]
If a $C_0$-group on a Banach space is analytic for positive time, then its
generator is bounded.
\end{lemma}

\begin{proof}
Fix $s>0$.  Analyticity implies
$T(s)X\subset\Dom(A)$ and $AT(s)\in\Lin(X)$.  Since $T(-s)$ is bounded and
$T(s)T(-s)=I$, every $x\in X$ belongs to $\Dom(A)$ and
\[
 A=(AT(s))T(-s)\in\Lin(X).
\]
\end{proof}

For $z_0\in\mathbb C$ and $0<\theta<\pi$, write
\[
 S_\theta(z_0)=\{z_0-w:w\neq0,\ |\arg w|<\theta\}.
\]
We use the standard meaning of a bounded $H^\infty$ calculus on this shifted
left sector.  Such a calculus for $A$ makes $B=z_0-A$ sectorial of angle at
most $\theta$.  If $\theta<\pi/2$, the sectorial generation theorem makes
$e^{-tB}$ analytic, and therefore makes
$T(t)=e^{tz_0}e^{-tB}$ analytic as well \cite{Haase}.

\begin{theorem}[Full resolution for the hyperbolic flow]\label{thm:main}
Let $\beta$ be any positive weight sequence satisfying the source
hypotheses, and suppose $T(t)=C_{\alpha_{\tanh t}}$ is a $C_0$-semigroup on
$\Hbeta$.  Its generator $A$ has the following properties.
\begin{enumerate}
\item $A$ is unbounded.
\item For no $z_0\in\mathbb C$ and no $0<\theta<\pi/2$ does $A$ admit a
bounded $H^\infty$ calculus on $S_\theta(z_0)$.
\end{enumerate}
\end{theorem}

\begin{proof}
Part (1) is Lemma~\ref{lem:unbounded}.  If the calculus in (2) existed, the
preceding sectorial generation argument would make $T$ analytic.  The group
extension and the analytic-group lemma would then make $A$ bounded,
contradicting (1).
\end{proof}

\begin{corollary}[Answers to the two readings]
The literal existence question has answer \emph{no}: the stated generator is
never bounded.  If ``bounded'' was intended to modify the composition
semigroup, in the source's sense that every member is bounded, then every
admissible choice of weights has the asserted failure of all proper-sector
calculi.  In particular the classical choice $\beta_n=1$ is an example.
\end{corollary}

This does not contradict the half-plane statement immediately preceding the
question: the boundary angle $\pi/2$ is allowed there, while the theorem
excludes exactly the strict range $0<\theta<\pi/2$.

\section{Verification, limitations, and novelty}

The accompanying verifier checks the flow and reflection identities and the
adjacent-monomial lower bound for constant, alternating, polynomial,
geometric, and pseudorandom positive weights.  The checks pass through
$n=500$; they are sanity checks for the exact algebraic proof, not a
substitute for it.

\textbf{Adversarial checks.}
No regularity of the ratios $\beta_{n+1}/\beta_n$ is assumed.  The proof uses
two adjacent normalized monomials precisely to prevent extreme ratios from
evading the estimate.  Polynomials belong to the generator domain by
\eqref{eq:generator}.  The proper-sector conclusion uses bounded
$H^\infty$ calculus in the standard sectorial sense; a half-plane calculus
at angle $\pi/2$ remains possible.  The group argument does not assume
uniform boundedness in time.

\textbf{Scope limitation.}
The theorem fully resolves the fixed hyperbolic-flow question, including its
wording ambiguity.  It does not classify $H^\infty$ calculi for arbitrary
noninvertible composition semiflows, nor does it claim that every
composition-semigroup generator is unbounded.

\textbf{Bounded novelty audit.}
The four run indexes and focused primary-source searches through
13 August 2026 were checked using arXiv:1708.02259 and combinations of
``weighted Hardy'', ``hyperbolic composition'', ``generator'',
``sectorial'', and ``$H^\infty$ calculus''.  No later paper explicitly
answering the displayed question was found.  Novelty is plausible, not
certified.

\textbf{Human-review recommendation:} send to an operator-semigroup or
composition-operator expert.  The highest-value review points are the
interpretation of ``bounded'' in the source wording and the standard
implication from a proper-sector bounded $H^\infty$ calculus to analyticity.

\begin{thebibliography}{9}
\bibitem{GY}
E.~A. Gallardo-Guti\'errez and D.~V. Yakubovich,
\emph{On generators of $C_0$-semigroups of composition operators},
Israel J. Math. \textbf{229} (2019), 487--500;
arXiv:1708.02259.

\bibitem{Haase}
M.~Haase,
\emph{The Functional Calculus for Sectorial Operators},
Operator Theory: Advances and Applications 169, Birkh\"auser, 2006.
\end{thebibliography}

\end{document}
